\section{Constante de Planck et droite d'action : dérivation déterminantale de l'échelle \texorpdfstring{\(10^{-34}\)}{décadique -34}}
\label{chap:escala-accion}
\index{Planck, constante de}\index{action!échelle de Planck}

La construction de l'échelle d'action commence par un invariant entier. La
forme normale de Smith du bloc dodécaphasique local produit un conoyau
d'ordre \(16\) ; sa norme sur la section constante \(\alphaHMT\), suivie d'une
unique action de \(\varphi\), place le résultat
\(\Sigma_H=\varphi\alphaHMT^{16}\) entre \(10^{-34}\) et \(10^{-33}\). De
cette manière, la puissance décimale caractéristique de l'action s'obtient avant
le choix des joules et des secondes. La branche angulaire fournit ensuite les deux
coordonnées orientées --- antérieure au retour et postérieure au retour ---, et la
carte métrologique les exprime finalement comme \(h\) et \(\hbar\).

\begin{exactbox}[title={Résultat directeur : du conoyau local à l'échelle d'action}]
La forme de Smith de \(H_4\), l'ordre de son conoyau et les inégalités qui
déterminent la décennie sont exacts. Le résultat complet est
\textsc{démontré avec une structure initiale explicite} : sont déclarés le
bloc local, la section invariante, la norme multiplicative, une unique
action de \(\varphi\) et le lecteur décimal de
\cref{chap:orden-determinantal-accion}. Les ablations identifient la fonction
de chaque règle et distinguent le conoyau local du conoyau global d'ordre
\(4096\).
\end{exactbox}

\subsection{Forme normale de Smith de \texorpdfstring{\(H_4\)}{H4} et transport de \texorpdfstring{\(\Sigma_H=\varphi\alpha_{\mathrm{HMT}}^{16}\)}{Sigma H} vers la droite d'action}

Pour l'une des trois orbites dodécaphasiques parallèles, les Parties~I--III obtiennent
\[
\operatorname{SNF}(H_4)=\operatorname{diag}(1,2,2,4),
\qquad
G_H:=\operatorname{coker}(H_4)
\cong\ZZ/2\ZZ\oplus\ZZ/2\ZZ\oplus\ZZ/4\ZZ,
\]
et, par conséquent, \(|G_H|=16\). La localité appartient au type du lecteur : le
conoyau de la transformation dodécaphasique complète est d'ordre
\(16^3=4096\), mais il enregistre simultanément trois copies directes et ne
remplace pas l'indice d'une orbite.

Soit \(s_\alpha(g)=\alphaHMT\) la section constante sur \(G_H\). Avec la
norme locale
\[
\mathcal N_{G_H}(s):=\prod_{g\in G_H}s(g)
\]
et une seule action de \(\varphi\) sur la base, on retrouve, avec une
étiquette locale facilitant les références de la Partie~V,
\begin{equation}
\label{eq:sigma-accion}
\Sigma_H
:=\varphi\,\mathcal N_{G_H}(s_\alpha)
=\varphi\alphaHMT^{16}.
\end{equation}
C'est la même section que celle définie dans \cref{eq:sigma-h-local}, et non une seconde
construction.

\begin{theorem}[Transfert de l'ordre déterminantal]
\label{thm:orden-accion}
Les primitives du lecteur déterminantal local étant fixées, la section de
\cref{eq:sigma-accion} satisfait
\[
\alphaHMT^{16}<10^{-34}
<\varphi\alphaHMT^{16}<10^{-33},
\qquad
\varepsilon_H:=\left\lfloor\log_{10}\Sigma_H\right\rfloor=-34.
\]
Si \(\mathcal L_S\) est une droite réelle d'action de base
\(\mathcal U_S\), l'affectation
\[
\tau_S:\RR_{>0}\longrightarrow\mathcal L_S,
\qquad x\longmapsto x\,\mathcal U_S,
\]
transporte cet ordre sans le modifier. En particulier, l'exposant décimal
des coordonnées d'action définies ci-dessous est \(-34\).
\end{theorem}

\begin{proof}
Les inégalités exactes et la valeur de l'ordre constituent le contenu de
\cref{thm:exponente-accion-local} ; elles s'y réduisent à des comparaisons entières à
partir du décimal rationnel construit de \(\alphaHMT\). L'application
\(\tau_S\) change le codomaine --- d'un scalaire positif vers une droite de
grandeurs à base déclarée ---, mais conserve la coordonnée scalaire. Par
conséquent, elle n'en modifie pas l'ordre décimal.
\end{proof}

L'évaluation héritée est
\[
\Sigma_H
=1.046243838104756580184229208303089\ldots\times10^{-34},
\]
et le nombre normalisé
\[
B_H:=10^{34}\Sigma_H
=1.046243838104756580184229208303089\ldots
\]
appartient à \([1,10)\). Le facteur \(10^{34}\) est reconstruit à partir de
l'ordre ; il n'est pas introduit comme seuil ou normalisation préalable.

\subsection{Nécessité du conoyau local}

Les quatre ablations exactes sont
\[
\begin{array}{c|c}
\text{règle évaluée}&\text{exposant décadique}\\
\hline
\alphaHMT^{16}\ \text{sans }\varphi&-35\\
\varphi\alphaHMT^{15}&-32\\
\varphi\alphaHMT^{17}&-37\\
\varphi\alphaHMT^{4096}&-8753.
\end{array}
\]
Ainsi, \(|G_H|=16\) fixe l'exposant de \(\alphaHMT\) une fois le
lecteur local choisi, et les inégalités fixent alors l'ordre \(-34\). Ces
contrôles excluent l'omission de \(\varphi\), les exposants voisins et le
remplacement du conoyau local par le conoyau global. Ils ne constituent ni une
classification de tous les lecteurs admissibles ni une preuve préalable d'unicité
de la section constante, de l'action unique de \(\varphi\) ou de la
carte décimale.

\subsection{Couplage ultérieur avec les coordonnées d'action}

Écrivons
\[
\varepsilon_H:=\left\lfloor\log_{10}\Sigma_H\right\rfloor=-34.
\]
Une fois construits par la branche angulaire le caractère antérieur au retour \(H_5\)
et la coordonnée reconstruite
\[
\etaRet=\frac{\Cstar}{2}-\alphaAngle,
\]
les deux reçoivent le même ordre dans \(\mathcal L_S\) :
\[
\hbar_{\HMT}^{\rm pre}
:=H_5\,10^{\varepsilon_H}\mathcal U_S,
\qquad
\hbar_{\HMT}^{\rm ret}
:=\etaRet\,10^{\varepsilon_H}\mathcal U_S.
\]
Numériquement,
\[
\hbar_{\HMT}^{\rm pre}
=1.054571817646154469625821829433059\ldots
\times10^{-34}\mathcal U_S,
\]
\[
\hbar_{\HMT}^{\rm ret}
=1.054571816915039061133690584724300\ldots
\times10^{-34}\mathcal U_S.
\]
La donnée \(H_5\) est antérieure au retour et ne s'identifie pas à \(\etaRet\).
Elle n'intervient pas non plus dans la démonstration de l'exposant \(-34\).

Les coordonnées non réduites correspondantes sont
\[
h_{\HMT}^{\rm pre}
=2\pi H_5\,10^{-34}\mathcal U_S
=6.626070149999987926001110349899474\ldots
\times10^{-34}\mathcal U_S,
\]
\[
h_{\HMT}^{\rm ret}
=2\pi\etaRet\,10^{-34}\mathcal U_S
=6.626070145406254333510749847592879\ldots
\times10^{-34}\mathcal U_S.
\]
La comparaison
\[
2\pi H_5-6.62607015
=-1.2073998889650100526\ldots\times10^{-14}
\]
est un contrôle externe ultérieur de la mantisse ; elle n'est ni une égalité avec la
valeur définitoire du SI ni une entrée du lecteur.

\begin{center}
\begin{tabularx}{0.96\textwidth}{@{}L{0.22\textwidth}L{0.36\textwidth}Y@{}}
\toprule
Objet & Coordenada obtenida & Fonction mathématique et métrologique\\
\midrule
Escala determinantal &
\(\Sigma_H=1.0462438381047565\ldots\times10^{-34}\) &
Fixe intérieurement l'ordre décimal de la droite d'action.\\
Action réduite préalable &
\(\hbar_{\rm HMT}^{\rm pre}=H_5\,10^{-34}\mathcal U_S\) &
Conserve la coordonnée angulaire antérieure au retour.\\
Action réduite postérieure &
\(\hbar_{\rm HMT}^{\rm ret}=\etaRet\,10^{-34}\mathcal U_S\) &
Donne la coordonnée après retour sur la même droite.\\
Action de cycle complet &
\(h_{\rm HMT}^{\rm pre}=6.6260701499999879\ldots\times10^{-34}\mathcal U_S\) &
Transporte l'action locale sur un tour de phase au moyen de \(h=2\pi\hbar\).\\
Carta SI &
\(h=6.62607015\times10^{-34}\ \mathrm{J\,s}\) &
Coordonnée définitoire externe utilisée pour la reconnaissance métrologique
\cite{BIPMSI}.\\
\bottomrule
\end{tabularx}
\end{center}

\subsection{\texorpdfstring{\(h\)}{h} et
\texorpdfstring{\(\hbar\)}{h barre} : action par cycle complet et par radian}
\label{sec:accion-vuelta-radian-rev6}
\index{action!par tour}
\index{action!par radian}

Les identités
\[
 E=h\nu=\hbar\omega,\qquad
 \omega=2\pi\nu,\qquad
 h=2\pi\hbar
\]
ne décrivent pas deux quanta. Ils décrivent une même action dans deux cartes. La
fréquence \(\nu\) compte les tours complets par unité de temps ; la
fréquence angulaire \(\omega\) compte la phase en radians par unité de temps.
En conséquence, \(h\) mesure l'action associée à la clôture d'un tour et
\(\hbar\) mesure l'action par unité de phase. La réduction n'est pas seulement une
abréviation : elle transporte l'holonomie globale vers le générateur local de sa
propagation.

La même distinction explique l'usage coordonné des degrés et des radians. Le cycle
angulaire complet de \(360^\circ\) rend visibles les partitions finies
\[
 12\cdot30^\circ=4\cdot90^\circ
 =3\cdot120^\circ=360^\circ
\]
et les retours de \(270^\circ\) et \(360^\circ\). Le radian est, en revanche, la coordonnée additive dans laquelle l’exponentielle et la dérivée s’écrivent sans insérer à chaque opération une constante de conversion. Les deux cartes sont reliées par
\[
 \mathbb R/\mathbb Z
 \xrightarrow{\;2\pi\;}
 \mathbb R/(2\pi\mathbb Z)
 \xrightarrow{\;180/\pi\;}
 \mathbb R/(360\mathbb Z).
\]
La composition conserve le point du tour : le degré exprime sa position dans le calendrier global et le radian linéarise son mouvement local.

Appliqué aux coordonnées construites ci-dessus,
\[
 h_{\rm HMT}^{\rm pre}
 =2\pi\,\hbar_{\rm HMT}^{\rm pre},\qquad
 h_{\rm HMT}^{\rm ret}
 =2\pi\,\hbar_{\rm HMT}^{\rm ret}.
\]
Le facteur \(2\pi\) ne corrige pas rétrospectivement une mantisse. Il parcourt avec l’action locale un tour complet. C’est pourquoi une amplitude \(\exp(\mathrm iS/\hbar)\), une relation de commutation et une quantification d’holonomie comparent une action à l’unité locale de phase, tandis que la fermeture globale s’exprime par \(h\).

\subsection{\(2\) branches causales et leur confluence}

La branche angulaire et la branche déterminantale partent de \(\alphaHMT\), mais aucune n’engendre rétrospectivement l’autre. La première se décompose comme suit
\[
\begin{gathered}
(\alphaHMT,\varphi;\,9,5,4,12,54,\text{ règle graduée})
\longrightarrow H_5,\\
\alphaHMT\longrightarrow A=1000\alphaHMT\,{}^\circ,
\qquad
(\alphaHMT,\pi)\longrightarrow D_A,\\
\mathcal F_\pi\longrightarrow\rho_\pi=169,\qquad
(\rho_\pi,D_A,\alphaHMT)\longrightarrow\mathcal C_\pi(\alphaHMT),\\
(\pi,\alphaHMT,H_5,\mathcal C_\pi(\alphaHMT))
\longrightarrow y_*\longrightarrow\Cstar,\\
(A,\Cstar)\longrightarrow q_\pm
\longrightarrow\bigl(S_{90}(q_\pm),S_{120}(q_\pm)\bigr)
\longrightarrow K_{6\to8}
\longrightarrow\Gamma_{6\to8}.
\end{gathered}
\]
L’identification ultérieure \(\Gamma_{6\to8}\dashrightarrow\gamma_{\rm BI}\) relève de l’interface physique. La branche déterminantale est
\[
(H_4,\alphaHMT,\varphi;\text{ règles du lecteur})
\longrightarrow G_H
\longrightarrow\Sigma_H
\longrightarrow\varepsilon_H=-34.
\]
Ce n’est qu’à la fin que \(\varepsilon_H\) se combine avec \(H_5\) et \(\etaRet\) pour donner des coordonnées dans \(\mathcal L_S\). Sont exclues aussi bien une flèche \(A\to H_5\) que toute sélection de \(C^*\) à partir du SI ou de \(\gamma_{\rm BI}\).

\subsection{Carte métrologique ultérieure et évaluation dimensionnelle des invariants engendrés}

Choisir une unité physique est une opération postérieure au théorème. Une carte métrologique est un isomorphisme de droites de grandeurs
\[
\iota_{\rm SI}:\mathcal L_S\longrightarrow\mathcal L_{\mathrm{J\,s}},
\qquad
\iota_{\rm SI}(\mathcal U_S)=\mathrm{J\,s}.
\]
L’unité appartient au codomaine de \(\iota_{\rm SI}\) ; l’exposant \(-34\) appartient au théorème arithmétique relatif. La juxtaposition d’un symbole d’unité ne transforme pas à elle seule un nombre sans dimension en une grandeur physique.

\begin{statusbox}[title={Hiérarchie probatoire de la constante de Planck}]
L’ordre \(10^{-34}\) provient du déterminant local et est fixé avant la métrologie. Les mantisses \(H_5\) et \(\etaRet\) proviennent de la branche angulaire et distinguent les orientations antérieure et postérieure au retour. L’identité \(h=2\pi\hbar\) transporte les deux coordonnées entre action par radian et action par cycle. L’application \(\iota_{\rm SI}\) fournit ensuite l’unité \(\mathrm{J\,s}\) et permet de comparer la sortie à la constante définitoire du SI. Cette succession maintient distinctes la dérivation interne, la réalisation dimensionnelle et la confrontation externe.
\end{statusbox}
